% !PRINTED 369
\weberpage{369}
%<<CONT
five complex formations of the type (3). But two such complex formations can differ only in this, that in place of $x_7$ there occurs each time another of the elements $x_3, x_4, x_5, x_6, x_7$, and all these possibilities must also occur, because otherwise we should obtain two different complexes which contain the same six elements of the complete system of the $x$, which is not possible, as we have proved. Accordingly we may assume that the $x_1, x_2, x_3, x_4, x_5, x_6$ occurring in the complex (3) $\beta$) are \gesperrt{any six elements whatever} of the complete system, which was to be proved.

\S.\,99.

\gesperrt{The Hesse-Cayley algorithm for the notation of the double tangents.}

Theorem 7. of the preceding section leads to a mode of notation for the double tangents by which a very perspicuous representation of all these relations is possible, and which has been developed by \gesperrt{Cayley} (following \gesperrt{Hesse}).

We take as a basis a complete seven-system

\[ x_1,\ x_2,\ x_3,\ x_4,\ x_5,\ x_6,\ x_7 \]

whose elements we denote simply by the figures 1, 2, 3, 4, 5, 6, 7. We separate off an arbitrary double tangent, say $x_1$, of this system, and form the complex which contains the remaining six:

\weq{1}{x_2 y_2,\ x_3 y_3,\ x_4 y_4,\ x_5 y_5,\ x_6 y_6,\ x_7 y_7.}

The double tangent $y_2$ is then completely determined by the chosen $x_1$ and by $x_2$, and can therefore be denoted by $[1\,2]$. In the same way we denote $y_3$ by $[1\,3]$, and so on. The double tangents $[1\,2], [1\,3], \dots, [1\,7]$ are hereby completely determined and different from one another. From this determination it is also evident what is to be understood in general by $[\mu\,\nu]$, when $\mu, \nu$ denote two different figures from the sequence 1 to 7.

It is now first to be shown that $[\mu\,\nu] = [\nu\,\mu]$. It suffices if we prove that $[1\,2] = [2\,1]$. For this we need only form the complex which contains $x_1, x_3, x_4, x_5, x_6, x_7$
%CONT>>
